\section{Attack Framework}
\label{sec:method}

\subsection{Threat Model and Access Spectrum}
\label{sec:threat-model}

We consider a graph $G=(V,E,X)$ with node features
$X\in\mathbb{R}^{|V|\times d}$, node labels $Y$, and a GNN classifier
$f_\theta$ trained on training mask
$\mathcal{D}_{\mathrm{tr}}\subset V$. A graph unlearning algorithm
$\mathcal{U}(f_\theta,\mathcal{D}_{\mathrm{tr}},S)\to f_{\theta'}$ takes a
deletion set $S\subset\mathcal{D}_{\mathrm{tr}}$ and returns an updated
model approximating retrain-from-scratch on
$\mathcal{D}_{\mathrm{tr}}\setminus S$.

\paragraph{Attacker capabilities.}
The attacker controls a deletion budget of
$k=\lceil r\cdot|\mathcal{D}_{\mathrm{tr}}|\rceil$ nodes, where $r$ is the
deletion ratio (we focus on $r=0.05$; \S\ref{app:ratio} reports the sweep).
The attacker selects $S\subset\mathcal{D}_{\mathrm{tr}}$ subject to
$|S|\le k$. Realistic deletion-API platforms motivate a
\emph{coalition-of-users} threat model: the attacker is a set of users who
jointly choose which records to forget. We do not assume the attacker can
modify training data or the unlearning algorithm itself.

\paragraph{Access spectrum.}
Selectors differ substantially in what information they need, and a
two-tier black/grey/white split obscures this. We use a finer, four-tier
spectrum keyed on \emph{resource cost}, not on which API endpoint is
exposed. We use L0 for random API-only deletion, L1 for structural
selectors available to a coalition with graph access, L2-\textsc{direct}
for model-coupled selectors using production weights, and
L2-\textsc{surrogate} for the shadow-model variant (Table~\ref{tab:access}
in Appendix~\ref{app:selectors}). The fingerprint axes
(\S\ref{sec:results-fingerprint}) report L1 (IM) and L2-\textsc{direct}
(TracIn); L2-\textsc{surrogate} is bounded above by the latter.

\paragraph{Attacker objective.}
Let $L(f;V_{\mathrm{te}})$ be downstream test F1. The attacker maximizes
the performance drop:
\begin{equation}
\label{eq:attacker-objective}
S^* = \argmax_{S\subset\mathcal{D}_{\mathrm{tr}},\,|S|\le k}\;
      L(f_\theta;V_{\mathrm{te}})
      - L\bigl(\mathcal{U}(f_\theta,\mathcal{D}_{\mathrm{tr}},S);V_{\mathrm{te}}\bigr).
\end{equation}
Because $\mathcal{U}$ is approximate, the attacker seeks deletion sets
that amplify approximation error beyond ordinary data removal. In
practice, selectors use public validation data as a proxy for
$V_{\mathrm{te}}$.

\subsection{Selection Strategies}
\label{sec:strategies}

We instantiate six strategies, tagged with their access tier from
Table~\ref{tab:access}.

\paragraph{Selectors.}
\textsc{Random} (L0) draws $k$ nodes uniformly without replacement.
\textsc{Degree} and \textsc{PageRank} (L1) choose high-centrality
training nodes. \textsc{IM} (L1) maximizes independent-cascade spread on
the graph, with the Monte-Carlo seed decoupled from GNN training seeds.
\textsc{TracIn} (L2-\textsc{direct}) ranks nodes by checkpoint-gradient
inner products against validation loss~\cite{pruthi2020estimating}.
\textsc{Hybrid} (L2) min-max normalizes and mixes the IM and TracIn
scores with $\alpha{=}0.5$. The fingerprint axes use IM and TracIn as
the structural-spread and gradient-signal endpoints; implementation
details and formulas are in Appendix~\ref{app:selectors}.

\subsection{Collateral Diagnostics}
\label{sec:diagnostics}

F1 drop alone conflates data removal, approximation drift, and
attack-specific signal. We therefore report the following diagnostics.

\paragraph{F1-shift decomposition.}
For any deletion set $S$, define
$\Delta F_{\text{total}}(S) = F_1(f_\theta;V_{\text{te}}) - F_1(\mathcal{U}(f_\theta,\cdot,S);V_{\text{te}})$
as the observed F1 \emph{drop} after unlearning (such that a positive
value indicates utility degradation). We decompose this response into
three interpretable terms. Let $R_5$ be a uniformly random set of five
training nodes, and let $R_r$ be a uniformly random set of
$\lceil r|\mathcal{D}_{\text{tr}}|\rceil$ training nodes:
\begin{align}
\Delta F_{\text{noise}}
&= \mathbb{E}_{R_5}
   [\Delta F_{\text{total}}(R_5)], \\
\Delta F_{\text{rand}}(r)
&= \mathbb{E}_{R_r}
   [\Delta F_{\text{total}}(R_r)], \\
\Delta F_{\text{volume}}(r)
&= \Delta F_{\text{rand}}(r)-\Delta F_{\text{noise}}, \\
\Delta F_{\text{attack}}(S)
&= \Delta F_{\text{total}}(S)-\Delta F_{\text{rand}}(r).
\end{align}
Thus, for an attack set $S$ of ratio $r$,
\begin{equation}
\label{eq:fshift}
\Delta F_{\text{total}}(S)
= \Delta F_{\text{noise}}
 + \Delta F_{\text{volume}}(r)
 + \Delta F_{\text{attack}}(S).
\end{equation}
$\Delta F_{\text{noise}}$ is a near-zero-volume noise floor: it measures
the intrinsic response of an unlearning method when only five random
training nodes are removed. We use ``noise floor'' in the empirical
diagnostic sense; it is not assumed to be zero-mean, and a nonzero value
is itself informative. $\Delta F_{\text{volume}}$ captures the additional
shift from increasing random deletion to the attack budget, while
$\Delta F_{\text{attack}}$ is the selector-specific excess over the
budget-matched random baseline. Only the final term is attack-specific;
we plot $\Delta F_{\text{attack}}$ in the fingerprint.

\paragraph{Retrain gap.} Let $f^{\mathrm{R}}$ be the model retrained from
scratch on $\mathcal{D}_{\mathrm{tr}}\setminus S$ and $f^{\mathrm{U}}$ be
the unlearned model. The retrain gap is
\[
\mathrm{Gap}(S)
= F_1(f^{\mathrm{R}};V_{\mathrm{te}})
 - F_1(f^{\mathrm{U}};V_{\mathrm{te}}).
\]
A large absolute gap indicates approximation drift relative to exact
retraining, not merely the informativeness of the deleted data.

\paragraph{Prediction shift.} Prediction shift measures retain-set
collateral perturbation relative to exact retraining. For retained nodes
$V_{\mathrm{ret}}$, we report the mean posterior shift
\[
\mathrm{Shift}(S)
= |V_{\mathrm{ret}}|^{-1}
\sum_{v\in V_{\mathrm{ret}}}
\|p^{\mathrm{U}}(v)-p^{\mathrm{R}}(v)\|_2,
\]
and the fraction of retained nodes whose predicted class differs between
$f^{\mathrm{U}}$ and $f^{\mathrm{R}}$.

\paragraph{Hop-distance decay.}
For each retained node $v\notin S$, let $h(v)=\min_{u\in S}d_G(v,u)$
be the shortest-path graph distance to the nearest deleted node. We
report \emph{flip rate per hop bucket}: the fraction of nodes in each
bucket whose predicted class differs between the unlearned model
$f^{\mathrm{U}}$ and the exact-retrain model $f^{\mathrm{R}}$ on
$\mathcal{D}_{\mathrm{tr}}\setminus S$:
\begin{equation}
\label{eq:hopflip}
\mathrm{Flip}(h) = \frac{1}{|B_h|}\sum_{v\in B_h}\mathbb{1}\bigl[f^{\mathrm{U}}(v)\neq f^{\mathrm{R}}(v)\bigr],
\qquad
B_h = \{v\notin S : h(v)\in\text{bucket } h\}.
\end{equation}
Using retrain as the oracle isolates approximation drift from data loss.
We bin by the GCN's effective receptive field rather than absolute hop
count: \textbf{within-RF} ($h\le L$), \textbf{boundary} ($h=L+1$), and
\textbf{outside-RF} ($h>L+1$). Sharp decay indicates localized drift;
flat curves indicate propagation beyond the nominal receptive field.
